\clearpage
\section{Problem Analysis}

\texttt{FSStatLib} computes a report from three input parameters:

\begin{itemize}
  \item \textbf{\texttt{D}}: the root directory to scan;
  \item \textbf{\texttt{MaxFS}}: the maximum file size used as threshold;
  \item \textbf{\texttt{NB}}: the number of regular size bands.
\end{itemize}

As a sequential computation, the problem consists of visiting each file,
reading its size, updating the total count, and incrementing the corresponding
histogram band.
The relevant part for this assignment is doing the same work
\textbf{asynchronously}, while still supporting \textbf{progress updates} and
\textbf{cancellation}.

\subsection{Filesystem Traversal}

The input directory may contain a large and irregular tree. A scanner must
therefore handle:

\begin{itemize}
  \item \textbf{recursion}: nested subdirectories must be visited;
  \item \textbf{scale}: the tree may contain many small files or many
  subdirectories;
  \item \textbf{robustness}: on real filesystems, entries may change while the
  scan is running, so transient access errors should not crash the program;
  \item \textbf{invalid or special entries}: invalid paths, unreadable
  directories, symbolic links, and special files must not crash or hang the
  program;
  \item \textbf{cycles}: repeated logical directories, for example through
  symbolic links, must not cause infinite traversal.
\end{itemize}

The scan is mainly \textbf{I/O-bound}. The expensive part is not the arithmetic
needed to compute a band index, but the number of filesystem operations needed
to list directories and inspect files. This makes responsiveness more important
than raw CPU throughput.

\subsection{Concurrency Requirements}

The asynchronous API must \textbf{return immediately} instead of waiting for
the whole directory tree to be processed. At the same time, the implementation
must coordinate work that may still be in progress. In particular, it must
define:

\begin{itemize}
  \item \textbf{scheduling}: how new directory work is started;
  \item \textbf{state ownership}: where partial counters are stored and
  updated;
  \item \textbf{completion detection}: how the implementation detects that no
  work is left;
  \item \textbf{progress reporting}: how updates are produced without flooding
  the CLI or GUI;
  \item \textbf{cancellation}: what happens when cancellation is requested
  while work is already in flight.
\end{itemize}

These requirements are solved differently in the three versions:

\begin{itemize}
  \item \textbf{Event-loop}: callbacks and pending asynchronous operations;
  \item \textbf{Reactive}: stream accumulation;
  \item \textbf{Virtual threads}: task coordination.
\end{itemize}

\subsection{Report Semantics}

The report contains two logical results:

\begin{itemize}
  \item \textbf{total file count}: the number of regular files found under
  \texttt{D}, including recursively visited subdirectories;
  \item \textbf{file-size histogram}: an array of \texttt{NB + 1} counters.
\end{itemize}

The first \texttt{NB} counters represent the bands obtained by splitting
\texttt{[0, MaxFS]}. The last counter is the \textbf{overflow band} and counts
files larger than \texttt{MaxFS}. Boundary values are handled
deterministically through a discrete partition of the integer size domain:
a file exactly equal to \texttt{MaxFS} belongs to the \textbf{last regular
band}, while a file larger than \texttt{MaxFS} belongs to the
\textbf{overflow band}.

The implementation counts \textbf{regular files} and skips non-regular
filesystem entries. This avoids treating directories, devices, or unsupported
links as ordinary files.

\subsection{Correctness Properties}

The expected behaviour can be summarized through these properties:

\begin{itemize}
  \item \textbf{Safety}: every counted file contributes to exactly one size
  band;
  \item \textbf{Completeness}: on a stable directory tree, every reachable
  regular file is counted once;
  \item \textbf{Consistency}: the three implementations produce the same report
  on the same stable directory tree;
  \item \textbf{Responsiveness}: progress updates and GUI interactions are not
  blocked by the traversal;
  \item \textbf{Termination}: completion, errors, and cancellation leave the CLI
  or GUI in a well-defined state.
\end{itemize}
